\documentclass[twocolumn]{aastex631}
% Auto-generated by jansky_research.fashienv._write_macros -- do not edit.
% Synthetic (feSyn*) and real (feReal*) namespaces are BOTH always emitted; the
% inactive namespace holds placeholders, so synthetic numbers can never masquerade
% under feReal* (an offline rebuild resets feReal* to placeholders by design).
\newcommand{\feSource}{FASHI DR2 (arXiv:2606.31539, CSTCloud share) x Tempel+2017 groups x Douglass+2023 voids (all holes); weights 1/(C*Vmax) from the DR2 catalogue (C >= 0.5)}
\newcommand{\feN}{110520}
\newcommand{\feSynVoidLogMStar}{--}
\newcommand{\feSynVoidAlpha}{--}
\newcommand{\feSynWallLogMStar}{--}
\newcommand{\feSynWallAlpha}{--}
\newcommand{\feSynGlobalLogMStar}{--}
\newcommand{\feSynGlobalAlpha}{--}
\newcommand{\feSynGroupLogMStar}{--}
\newcommand{\feSynFieldLogMStar}{--}
\newcommand{\feSynVoidKneeOffset}{--}
\newcommand{\feSynVoidKneeErr}{--}
\newcommand{\feSynVoidKneeSigma}{--}
\newcommand{\feSynVoidLogMStarErr}{--}
\newcommand{\feSynWallLogMStarErr}{--}
\newcommand{\feSynNWall}{--}
\newcommand{\feSynNField}{--}
\newcommand{\feSynGroupKneeOffset}{--}
\newcommand{\feSynGroupKneeErr}{--}
\newcommand{\feSynGroupKneeSigma}{--}
\newcommand{\feSynNInVoid}{--}
\newcommand{\feSynNGroupMembers}{--}
\newcommand{\feRealVoidLogMStar}{9.805}
\newcommand{\feRealVoidAlpha}{-1.266}
\newcommand{\feRealWallLogMStar}{9.976}
\newcommand{\feRealWallAlpha}{-1.429}
\newcommand{\feRealGlobalLogMStar}{9.907}
\newcommand{\feRealGlobalAlpha}{-1.328}
\newcommand{\feRealGroupLogMStar}{10.041}
\newcommand{\feRealFieldLogMStar}{9.841}
\newcommand{\feRealVoidKneeOffset}{-0.171}
\newcommand{\feRealVoidKneeErr}{0.047}
\newcommand{\feRealVoidKneeSigma}{3.610}
\newcommand{\feRealVoidLogMStarErr}{0.035}
\newcommand{\feRealWallLogMStarErr}{0.032}
\newcommand{\feRealNWall}{38939}
\newcommand{\feRealNField}{37210}
\newcommand{\feRealGroupKneeOffset}{0.200}
\newcommand{\feRealGroupKneeErr}{0.043}
\newcommand{\feRealGroupKneeSigma}{4.620}
\newcommand{\feRealNInVoid}{16954}
\newcommand{\feRealNGroupMembers}{18683}
\newcommand{\feRealNDRTwo}{156411}
\newcommand{\feRealNZNonpos}{286}
\newcommand{\feRealCMin}{0.5}
\newcommand{\feRealWallPct}{35.2}
\newcommand{\feRealWallMinusGlobal}{0.069}
\newcommand{\feRealGlobalLogMStarErr}{0.021}
\newcommand{\feRealGlobalAlphaErr}{0.032}
\newcommand{\feRealEdsVoidKneeOffset}{-0.076}
\newcommand{\feRealEdsVoidKneeSigma}{1.600}
\newcommand{\feRealOptAVoidKneeOffset}{-0.237}
\newcommand{\feRealOptAVoidKneeErr}{0.083}
\newcommand{\feRealOptAGroupKneeOffset}{0.289}
\newcommand{\feRealOptAGlobalAlpha}{-1.780}
\newcommand{\feRealDROneN}{41741}
\newcommand{\feRealDROneVoidKneeOffset}{-0.157}
\newcommand{\feRealDROneVoidKneeErr}{0.085}
\newcommand{\feRealDROneVoidKneeSigma}{1.850}
\newcommand{\feRealDROneGroupKneeOffset}{0.261}
\newcommand{\feRealDROneGlobalAlpha}{-1.732}
\newcommand{\feRealMatchedPct}{94.7}
\newcommand{\feRealNDropped}{142}
\newcommand{\feRealJkErr}{0.035}
\newcommand{\feRealJkNOcc}{550}
\newcommand{\feRealWeightEffect}{0.066}
\newcommand{\feRealWeightEffectErr}{0.020}
\newcommand{\feRealOptASameVoidKneeOffset}{-0.237}
\newcommand{\feRealOptASameGroupKneeOffset}{0.322}
\newcommand{\feRealMaxSphereVoidKneeOffset}{-0.163}
\newcommand{\feRealMaxSphereNInVoid}{6200}
\newcommand{\feRealNullReps}{1000}
\newcommand{\feRealNullMean}{-0.102}
\newcommand{\feRealNullStd}{0.022}
\newcommand{\feRealNullNLe}{0}
\newcommand{\feRealNullExcess}{-0.069}
\newcommand{\feRealNullNInVoid}{26708}
\newcommand{\feRealNullOverlap}{32}
\newcommand{\feRealNullCReps}{1000}
\newcommand{\feRealNullCMean}{-0.107}
\newcommand{\feRealNullCStd}{0.021}
\newcommand{\feRealNullCNLe}{0}
\newcommand{\feRealNullCExcess}{-0.064}
\newcommand{\feRealNullCNInVoid}{28154}
\newcommand{\feRealNullCNBins}{14}
\newcommand{\feRealNullCOverlap}{33}
\newcommand{\feRealNullCIntercept}{-0.265}
\newcommand{\feRealNullCOptAMean}{-0.063}
\newcommand{\feRealNullCSlope}{0.477}
\newcommand{\feRealNullIntercept}{-0.275}
\newcommand{\feRealNullSlope}{0.529}
\newcommand{\feRealNullCOptAStd}{0.048}
\newcommand{\feRealNullCOptAMin}{-0.187}
\newcommand{\feRealNullCOverlapLo}{0.285}
\newcommand{\feRealNullCOverlapHi}{0.371}
\newcommand{\feRealNullCNInVoidMin}{25050}
\newcommand{\feRealNullCNInVoidMax}{31423}
\newcommand{\feRealNullCRtwoOverlap}{0.083}
\newcommand{\feRealNullCRtwoOccZ}{0.284}
\newcommand{\feRealNullCRtwoAll}{0.285}
\newcommand{\feRealNullCOverlapCoefCtl}{0.046}
\newcommand{\feRealNullCOverlapChange}{0.041}
\newcommand{\feRealOptASameGlobalAlpha}{-1.729}
\newcommand{\feRealOptASameGlobalRedChi}{190}
\newcommand{\feRealGroupNullNInGroup}{1909}
\newcommand{\feRealGroupNInGroupReal}{18683}
\newcommand{\feRealGroupNullMedZ}{0.022}
\newcommand{\feRealGroupMedZReal}{0.031}
\newcommand{\feRealGroupNullOverlap}{44}
\newcommand{\feRealNoOverlapUnplaced}{688}
\newcommand{\feRealNoOverlapNVoids}{1163}
\newcommand{\feRealEdsPercentile}{92.7}
\newcommand{\feRealGroupNullReps}{200}
\newcommand{\feRealGroupNullMean}{0.127}
\newcommand{\feRealGroupNullStd}{0.059}
\newcommand{\feRealGroupNullNGe}{23}
\newcommand{\feRealGroupNullExcess}{0.073}
\newcommand{\feRealCommonKneeOffset}{-0.145}
\newcommand{\feRealCommonKneeErr}{0.047}
\newcommand{\feRealCommonNBins}{14}
\newcommand{\feRealVoidNBins}{14}
\newcommand{\feRealWallNBins}{18}
\newcommand{\feRealDROneMaxVoidKneeOffset}{-0.259}
\newcommand{\feRealGlobalRedChi}{41}
\newcommand{\feRealVoidRedChi}{19}
\newcommand{\feRealWallRedChi}{29}
\newcommand{\feRealDistRatio}{1.025}
\newcommand{\feRealShufVoidSurveyMean}{-0.078}
\newcommand{\feRealShufVoidSurveyStd}{0.015}
\newcommand{\feRealShufVoidSurveyExcess}{-0.093}
\newcommand{\feRealShufVoidSurveySigma}{1.9}
\newcommand{\feRealShufVoidSurveyNReach}{0}
\newcommand{\feRealShufZVoidSurveyExcess}{-0.100}
\newcommand{\feRealShufZVoidSurveySigma}{2.0}
\newcommand{\feRealShufTVoidSurveyExcess}{-0.091}
\newcommand{\feRealShufTVoidSurveySigma}{1.9}
\newcommand{\feRealShufVoidEnvMean}{0.015}
\newcommand{\feRealShufVoidEnvStd}{0.016}
\newcommand{\feRealShufVoidEnvExcess}{-0.103}
\newcommand{\feRealShufVoidEnvSigma}{2.5}
\newcommand{\feRealShufVoidEnvNReach}{0}
\newcommand{\feRealShufZVoidEnvExcess}{-0.110}
\newcommand{\feRealShufZVoidEnvSigma}{2.6}
\newcommand{\feRealShufTVoidEnvExcess}{-0.103}
\newcommand{\feRealShufTVoidEnvSigma}{2.5}
\newcommand{\feRealBlendVoidEnvAll}{-0.088}
\newcommand{\feRealBlendVoidEnvAllErr}{0.038}
\newcommand{\feRealBlendVoidSameGroupPct}{98.5}
\newcommand{\feRealBlendVoidWfiftyFlagIn}{6.8}
\newcommand{\feRealBlendVoidWfiftyFlagOut}{10.8}
\newcommand{\feRealBlendVoidWfiftyEnvClean}{-0.078}
\newcommand{\feRealBlendVoidWfiftyEnvCleanErr}{0.036}
\newcommand{\feRealBlendVoidWfiftyShufExcess}{-0.092}
\newcommand{\feRealBlendVoidWfiftyShufSigma}{2.3}
\newcommand{\feRealBlendVoidWfiftyLogMFlag}{9.88}
\newcommand{\feRealBlendVoidWfiftyLogMClean}{9.48}
\newcommand{\feRealBlendVoidFixedFlagIn}{7.5}
\newcommand{\feRealBlendVoidFixedFlagOut}{12.2}
\newcommand{\feRealBlendVoidFixedEnvClean}{-0.080}
\newcommand{\feRealBlendVoidFixedEnvCleanErr}{0.036}
\newcommand{\feRealBlendVoidFixedShufExcess}{-0.095}
\newcommand{\feRealBlendVoidFixedShufSigma}{2.4}
\newcommand{\feRealBlendVoidFixedLogMFlag}{9.86}
\newcommand{\feRealBlendVoidFixedLogMClean}{9.48}
\newcommand{\feRealBlendVoidLwInInnerDiff}{-0.007}
\newcommand{\feRealBlendVoidLwInInnerSe}{0.007}
\newcommand{\feRealBlendVoidLwInInnerSigma}{-0.9}
\newcommand{\feRealBlendVoidLwInInnerN}{1417}
\newcommand{\feRealBlendVoidLwInRingDiff}{0.015}
\newcommand{\feRealBlendVoidLwInRingSe}{0.007}
\newcommand{\feRealBlendVoidLwInRingSigma}{2.1}
\newcommand{\feRealBlendVoidLwInRingN}{1563}
\newcommand{\feRealBlendVoidLwInRoneDiff}{-0.006}
\newcommand{\feRealBlendVoidLwInRtwoDiff}{-0.002}
\newcommand{\feRealBlendVoidLwInRthreeDiff}{0.004}
\newcommand{\feRealBlendVoidLwOutInnerDiff}{-0.003}
\newcommand{\feRealBlendVoidLwOutInnerSe}{0.004}
\newcommand{\feRealBlendVoidLwOutInnerSigma}{-0.6}
\newcommand{\feRealBlendVoidLwOutInnerN}{5321}
\newcommand{\feRealBlendVoidLwOutRingDiff}{0.009}
\newcommand{\feRealBlendVoidLwOutRingSe}{0.004}
\newcommand{\feRealBlendVoidLwOutRingSigma}{2.2}
\newcommand{\feRealBlendVoidLwOutRingN}{5554}
\newcommand{\feRealBlendVoidLwOutRoneDiff}{-0.001}
\newcommand{\feRealBlendVoidLwOutRtwoDiff}{-0.001}
\newcommand{\feRealBlendVoidLwOutRthreeDiff}{0.007}
\newcommand{\feRealBlendVoidInnerN}{1443}
\newcommand{\feRealBlendVoidInnerLogM}{9.86}
\newcommand{\feRealBlendVoidRingN}{1578}
\newcommand{\feRealBlendVoidRingLogM}{9.59}
\newcommand{\feRealBlendVoidIsolatedN}{14323}
\newcommand{\feRealBlendVoidIsolatedLogM}{9.46}
\newcommand{\feRealBlendVoidRRFlag}{--}
\newcommand{\feRealBlendVoidRRClean}{--}
\newcommand{\feRealStrictVoidSurveyOffset}{-0.172}
\newcommand{\feRealStrictVoidEnvOffset}{-0.088}
\newcommand{\feRealShufGroupSurveyMean}{0.077}
\newcommand{\feRealShufGroupSurveyStd}{0.013}
\newcommand{\feRealShufGroupSurveyExcess}{0.123}
\newcommand{\feRealShufGroupSurveySigma}{2.7}
\newcommand{\feRealShufGroupSurveyNReach}{0}
\newcommand{\feRealShufZGroupSurveyExcess}{0.133}
\newcommand{\feRealShufZGroupSurveySigma}{2.9}
\newcommand{\feRealShufTGroupSurveyExcess}{0.123}
\newcommand{\feRealShufTGroupSurveySigma}{2.7}
\newcommand{\feRealShufGroupEnvMean}{-0.005}
\newcommand{\feRealShufGroupEnvStd}{0.012}
\newcommand{\feRealShufGroupEnvExcess}{0.128}
\newcommand{\feRealShufGroupEnvSigma}{3.5}
\newcommand{\feRealShufGroupEnvNReach}{0}
\newcommand{\feRealShufZGroupEnvExcess}{0.139}
\newcommand{\feRealShufZGroupEnvSigma}{3.8}
\newcommand{\feRealShufTGroupEnvExcess}{0.129}
\newcommand{\feRealShufTGroupEnvSigma}{3.6}
\newcommand{\feRealBlendGroupEnvAll}{0.123}
\newcommand{\feRealBlendGroupEnvAllErr}{0.034}
\newcommand{\feRealBlendGroupSameGroupPct}{99.4}
\newcommand{\feRealBlendGroupWfiftyFlagIn}{25.5}
\newcommand{\feRealBlendGroupWfiftyFlagOut}{1.6}
\newcommand{\feRealBlendGroupWfiftyEnvClean}{0.038}
\newcommand{\feRealBlendGroupWfiftyEnvCleanErr}{0.035}
\newcommand{\feRealBlendGroupWfiftyShufExcess}{0.062}
\newcommand{\feRealBlendGroupWfiftyShufSigma}{1.6}
\newcommand{\feRealBlendGroupWfiftyLogMFlag}{9.92}
\newcommand{\feRealBlendGroupWfiftyLogMClean}{9.47}
\newcommand{\feRealBlendGroupFixedFlagIn}{29.1}
\newcommand{\feRealBlendGroupFixedFlagOut}{1.6}
\newcommand{\feRealBlendGroupFixedEnvClean}{0.044}
\newcommand{\feRealBlendGroupFixedEnvCleanErr}{0.035}
\newcommand{\feRealBlendGroupFixedShufExcess}{0.076}
\newcommand{\feRealBlendGroupFixedShufSigma}{2.0}
\newcommand{\feRealBlendGroupFixedLogMFlag}{9.89}
\newcommand{\feRealBlendGroupFixedLogMClean}{9.46}
\newcommand{\feRealBlendGroupLwInInnerDiff}{-0.006}
\newcommand{\feRealBlendGroupLwInInnerSe}{0.006}
\newcommand{\feRealBlendGroupLwInInnerSigma}{-1.2}
\newcommand{\feRealBlendGroupLwInInnerN}{5952}
\newcommand{\feRealBlendGroupLwInRingDiff}{0.010}
\newcommand{\feRealBlendGroupLwInRingSe}{0.005}
\newcommand{\feRealBlendGroupLwInRingSigma}{1.9}
\newcommand{\feRealBlendGroupLwInRingN}{4398}
\newcommand{\feRealBlendGroupLwInRoneDiff}{-0.008}
\newcommand{\feRealBlendGroupLwInRtwoDiff}{-0.005}
\newcommand{\feRealBlendGroupLwInRthreeDiff}{0.001}
\newcommand{\feRealBlendGroupLwOutInnerDiff}{-0.027}
\newcommand{\feRealBlendGroupLwOutInnerSe}{0.009}
\newcommand{\feRealBlendGroupLwOutInnerSigma}{-2.9}
\newcommand{\feRealBlendGroupLwOutInnerN}{740}
\newcommand{\feRealBlendGroupLwOutRingDiff}{-0.001}
\newcommand{\feRealBlendGroupLwOutRingSe}{0.005}
\newcommand{\feRealBlendGroupLwOutRingSigma}{-0.1}
\newcommand{\feRealBlendGroupLwOutRingN}{2714}
\newcommand{\feRealBlendGroupLwOutRoneDiff}{-0.011}
\newcommand{\feRealBlendGroupLwOutRtwoDiff}{-0.027}
\newcommand{\feRealBlendGroupLwOutRthreeDiff}{0.000}
\newcommand{\feRealBlendGroupInnerN}{6051}
\newcommand{\feRealBlendGroupInnerLogM}{9.89}
\newcommand{\feRealBlendGroupRingN}{4419}
\newcommand{\feRealBlendGroupRingLogM}{9.60}
\newcommand{\feRealBlendGroupIsolatedN}{8662}
\newcommand{\feRealBlendGroupIsolatedLogM}{9.32}
\newcommand{\feRealBlendGroupRRFlag}{0.35}
\newcommand{\feRealBlendGroupRRClean}{0.52}
\newcommand{\feRealStrictGroupSurveyOffset}{0.200}
\newcommand{\feRealStrictGroupEnvOffset}{0.124}
\newcommand{\feRealBlendNPool}{55893}
\newcommand{\feRealStrictNPool}{55651}
\newcommand{\feRealStrictMinPerCell}{5}
\newcommand{\feRealShufReps}{1000}
\newcommand{\feRealStrictMaxShift}{0.001}
\newcommand{\feRealEnvVoidShareFirst}{46}
\newcommand{\feRealEnvVoidShareAll}{61}
\newcommand{\feRealEnvGroupShellFirst}{6}
\newcommand{\feRealEnvGroupShellLast}{1}
\newcommand{\feRealEnvVoidQuadSigma}{1.1}
\newcommand{\feRealEnvGroupQuadSigma}{2.2}
\newcommand{\feRealNullCExcessSigmaQuad}{1.2}
\newcommand{\feRealGroupNullExcessSigmaQuad}{1.0}
\newcommand{\feRealVoidBaseRatePct}{30}
\newcommand{\feRealVoidBracketLo}{0.05}
\newcommand{\feRealVoidBracketHi}{0.10}
\newcommand{\feRealEnvGroupNullMeanSe}{0.004}
\newcommand{\feRealEnvGroupNullPTwo}{0.05}
\newcommand{\feRealEnvVoidNullShift}{0.067}
\newcommand{\feRealEnvGroupNullShift}{-0.147}
\newcommand{\feRealEnvVoidMeasShift}{0.083}
\newcommand{\feRealEnvGroupMeasShift}{-0.077}
\newcommand{\feRealEnvNullCorrOccSurvey}{-0.45}
\newcommand{\feRealEnvNullCorrOccEnv}{-0.26}
\newcommand{\feRealEnvFrameRatio}{1.052}
\newcommand{\feRealEnvFrameBelowPct}{10}
\newcommand{\feRealEnvFrameMaxShift}{0.005}
\newcommand{\feRealEnvVoidOffsetNoFrame}{-0.093}
\newcommand{\feRealEnvGroupOffsetNoFrame}{0.123}
\newcommand{\feRealEnvVoidOffset}{-0.088}
\newcommand{\feRealEnvVoidErr}{0.038}
\newcommand{\feRealEnvGroupOffset}{0.123}
\newcommand{\feRealEnvGroupErr}{0.034}
\newcommand{\feRealEnvNullReps}{200}
\newcommand{\feRealEnvNullMean}{-0.040}
\newcommand{\feRealEnvNullStd}{0.020}
\newcommand{\feRealEnvNullNLe}{2}
\newcommand{\feRealEnvNullExcess}{-0.048}
\newcommand{\feRealEnvNullExcessSigma}{1.3}
\newcommand{\feRealEnvGroupNullReps}{200}
\newcommand{\feRealEnvGroupNullMean}{-0.020}
\newcommand{\feRealEnvGroupNullStd}{0.056}
\newcommand{\feRealEnvGroupNullNGe}{4}
\newcommand{\feRealEnvGroupNullP}{0.02}
\newcommand{\feRealEnvGroupNullExcessSigma}{4.2}
\newcommand{\feRealEnvGroupNullExcess}{0.143}
\newcommand{\feRealEnvVoidShellLo}{46}
\newcommand{\feRealEnvVoidShellHi}{67}
\newcommand{\feRealEnvGroupShellLo}{1}
\newcommand{\feRealEnvGroupShellHi}{6}
\newcommand{\feRealNullCOccCorr}{-0.49}
\newcommand{\feRealNullCOccSlope}{-0.011}
\newcommand{\feRealNullCExcessSigmaNull}{3.0}
\newcommand{\feRealNullCExcessSigmaFit}{1.4}
\newcommand{\feRealNullCExcessSigmaJk}{1.8}
\newcommand{\feRealNullCJkErr}{0.040}
\newcommand{\feRealNullCBiasPct}{63}
\newcommand{\feRealWeightEffectSigma}{3.3}
\newcommand{\feRealOptANullExcess}{-0.174}
\newcommand{\feRealEdsChangedLowZPct}{7}
\newcommand{\feRealEdsChangedHighZPct}{126}

\newcommand{\code}[1]{\texttt{#1}}
\shorttitle{Random-placement nulls for environment-split HIMFs}
\shortauthors{Barbere}
\begin{document}
\title{Environment-Split HI Mass Functions in FASHI DR2 and Random-Placement Null Tests}
\author[0009-0008-3289-4447]{Joseph Barbere}
\affiliation{Independent researcher}

\begin{abstract}
We split the HI mass function of FASHI DR2 by environment and compare the knee offsets with two
null tests: void- and group-shaped volumes placed at random, and environment labels shuffled
among galaxies at fixed redshift with the real number of members. With the survey-wide
$1/(C\,V_{\rm max})$ weighting the void$-$wall and group$-$field offsets are
$\feRealVoidKneeOffset$ and $+\feRealGroupKneeOffset$~dex; the random placements give
$\feRealNullCMean$ and $+\feRealGroupNullMean$~dex and the shuffles $\feRealShufVoidSurveyMean$
and $+\feRealShufGroupSurveyMean$~dex, so the two nulls disagree about the baseline. The void
offset exceeds the shuffle null by \feRealShufVoidSurveySigma--\feRealShufVoidEnvSigma$\sigma$
under two weightings and is unchanged when beam-confused sources are removed, but exceeds the
random placements by only \feRealEnvVoidQuadSigma--\feRealNullCExcessSigmaQuad$\sigma$. A quarter of
group members have an optical neighbour inside the FAST beam, against \feRealBlendGroupWfiftyConfOut\% of field galaxies, and
without them most of the group offset disappears. We report the void offset as bracketed by the
two nulls and do not attribute the group offset to environment.
\end{abstract}

\keywords{HI line emission (690), Galaxy environments (2029), Large-scale structure of the
universe (902), Luminosity function (942)}

\section{Introduction}
Void galaxies are expected to have a lower HI mass function (HIMF) knee than wall galaxies, and
group galaxies to lose gas to stripping \citep{solanes2001,jones2018}; \citet{moorman2014}, with an
estimator that does not assume uniform density, found the ALFALFA void knee $\sim$0.14~dex lower.
The FAST All Sky HI Survey \citep[FASHI;][]{fashi_dr1,fashi_dr2} now provides \feRealNDRTwo\
sources with a completeness $C$ and accessible volume $V_{\rm max}$ for each; the one environment
study on FASHI data used 230 group galaxies \citep{fashi_groups}. Environment splits of a
flux-limited $1/V_{\rm max}$ HIMF also select in distance and structure, so we compare them with
two null tests that fail in different ways, and test the one systematic that follows density
directly, blending in FAST's beam. An earlier version of this analysis, on DR1 with a single
flux-limit weighting, reported a void offset of $\feRealDROneMaxVoidKneeOffset$~dex without a
null test.

\section{Data and Method}
\label{sec:data}
We keep the \feN\ sources of Table~2 of \citet{fashi_dr2} with $C\geq\feRealCMin$, the sample of
the survey's own HIMF; weights $1/(C_i V_i)$ reproduce the published global HIMF
($\log M^*=\feRealGlobalLogMStar$, $\alpha=\feRealGlobalAlpha$, against 9.89 and $-1.31$).
Schechter fits are formally poor (reduced $\chi^2$ up to \feRealGlobalRedChi); their errors are
scaled accordingly. Void members lie inside any hole of a Douglass+2023 VoidFinder void
\citep{douglass2023}, and walls are the rest of a padded box around the voids. Group members lie
within the projected $R_{200}$ and 1000~km\,s$^{-1}$ of a Tempel+2017 group \citep{tempel2017}.
Besides the survey-wide weighting we restrict each galaxy's $V_{\rm max}$ to its environment,
$C_i V_i\,g_E(d_{\max,i})$, where $g_E(D)$ is the share of the volume within $D$ occupied by
environment $E$, from $2\times10^6$ randoms classified by the same rules.

The first null moves each void rigidly to a random position with all its holes in the SDSS
footprint, and each group to a random footprint position with its velocity and $R_{200}$
(\feRealNullCReps\ and \feRealGroupNullReps\ placements). The second gives the environment label
to the real number of members in each $\Delta z=0.0025$ bin, drawn at random from the sample
(\feRealShufReps\ shuffles); it matches the occupancy and redshift distribution exactly but has no
spatial coherence. For blending, a source is confused when two or more Tempel SDSS galaxies
($r<17.77$), one of them its counterpart, lie within one FAST beam (2.9~arcmin) and $W_{50}/2+100$
km\,s$^{-1}$ of it; we use the \feRealBlendNSdss\ sources inside the SDSS footprint. Each excess
over a null is divided by the measured offset's fit error and the null spread in quadrature.

\section{Results}
\label{sec:results}
The measured offsets are $\feRealVoidKneeOffset\pm\feRealVoidKneeErr$ (void) and
$+\feRealGroupKneeOffset\pm\feRealGroupKneeErr$~dex (group) with the survey-wide weighting, and
$\feRealEnvVoidOffset\pm\feRealEnvVoidErr$ and $+\feRealEnvGroupOffset\pm\feRealEnvGroupErr$~dex
with the restricted one. Against the shuffles the void excess is $\feRealShufVoidSurveyExcess$ and
$\feRealShufVoidEnvExcess$~dex (\feRealShufVoidSurveySigma\ and \feRealShufVoidEnvSigma$\sigma$) and
the group excess $+\feRealShufGroupSurveyExcess$ and $+\feRealShufGroupEnvExcess$~dex
(\feRealShufGroupSurveySigma\ and \feRealShufGroupEnvSigma$\sigma$); no shuffle reaches either.
Against the random placements the void excess is $\feRealNullCExcess$ and
$\feRealEnvNullExcess$~dex (\feRealNullCExcessSigmaQuad\ and \feRealEnvVoidQuadSigma$\sigma$), and the group excess
$+\feRealGroupNullExcess$ and $+\feRealEnvGroupNullExcess$~dex (\feRealGroupNullExcessSigmaQuad\ and \feRealEnvGroupQuadSigma$\sigma$; Figure~\ref{fig:nulls}).

The two nulls differ in what they hold fixed. The random void placements contain
\feRealNullCNInVoidMin--\feRealNullCNInVoidMax\ galaxies against \feRealNInVoid\ in the real voids,
and the random group regions a median \feRealGroupNullNInGroup\ against \feRealGroupNInGroupReal;
their offsets correlate with occupancy and redshift, which the shuffles match. The shuffles lack
the spatial coherence of a real volume. Both include real members, at \feRealVoidBaseRatePct\% (void shuffles) and a median
\feRealNullCOverlap\% (void placements), which pulls each null towards the measurement.

Of the group members, \feRealBlendGroupWfiftyConfIn\% are confused, against
\feRealBlendGroupWfiftyConfOut\% of field galaxies; for voids and walls the fractions are
\feRealBlendVoidWfiftyConfIn\% and \feRealBlendVoidWfiftyConfOut\%. Removing confused sources
changes the restricted group offset from $+\feRealBlendGroupEnvAll$ to
$+\feRealBlendGroupWfiftyEnvClean$~dex, leaving $+\feRealBlendGroupWfiftyShufExcess$~dex
(\feRealBlendGroupWfiftyShufSigma$\sigma$) over the shuffle null; the void offset moves from
$\feRealBlendVoidEnvAll$ to $\feRealBlendVoidWfiftyEnvClean$~dex and keeps a
\feRealBlendVoidWfiftyShufSigma$\sigma$ excess. A fixed $\pm300$~km\,s$^{-1}$ window gives
$+\feRealBlendGroupFixedEnvClean$ and $\feRealBlendVoidFixedEnvClean$~dex. The flag selects in
mass under either window: confused group members have a median $\log M=\feRealBlendGroupWfiftyLogMConf$
against \feRealBlendGroupWfiftyLogMClean\ for the rest.

\begin{figure}
\plotone{figures/fashienv_nulls.pdf}
\caption{Knee offsets from \feRealEnvNullReps\ random placements of the voids (left) and groups
(right) under the survey-wide (grey) and environment-restricted (blue) weightings; shaded bands
are the label-shuffle nulls (mean $\pm1\sigma$) and dashed lines the measured offsets.}
\label{fig:nulls}
\end{figure}

\section{Discussion}
\label{sec:discussion}
The void offset is not explained by the redshift and occupancy of the void sample, and it does
not depend on beam confusion, but volumes of the same shapes placed at random reproduce most of
it. Placement preserves spatial structure and the shuffle does not; neither holds both occupancy
and structure fixed, and a void is defined by being emptier than any randomly placed volume, so
no placement can. We therefore report the void offset as \feRealVoidBracketLo--\feRealVoidBracketHi~dex beyond the two nulls
rather than as a detection.

The group offset is carried largely by the quarter of members that have a neighbour inside the
beam. Removing them also removes massive galaxies, whose neighbours are more often bright enough
for the optical catalogue, so the test cannot separate blending, which raises the HI mass, from
this selection. Interferometric HI for a subset of group members, or a forward model of confusion
in FAST's beam, would decide it; until then the group offset is not attributed to environment.
For environment-split $1/V_{\rm max}$ HIMFs from single-dish surveys we recommend both nulls and a
confusion test, with each null's occupancy reported.

Reproducibility: \code{jansky\_research.fashienv}; every number derived from the data is
pipeline-generated (\code{results/fashienv\_metrics.json} $\to$ macros); survey constants and
values quoted from other papers are typed. Data: FASHI DR2 Table~2 (public CSTCloud share linked
from the FASHI release page), FASHI DR1 for the comparison (VizieR \code{J/other/SCPMA/67.19511}),
Tempel+2017 (\code{J/A+A/602/A100}) and Douglass+2023 (\code{J/ApJS/265/7}, table1 and table2) via
VizieR; no authentication is needed.

\section*{AI use disclosure}
This paper was developed collaboratively by the author with Anthropic's Claude
\citep{anthropicclaude} (models Claude Opus~4.x and 5.5, accessed 2026). Claude contributed to the design,
implementation, and testing of the \code{fashienv} module and the drafting of this manuscript.
The author directed the work, independently reviewed and verified all code, results, and
citations, and takes full responsibility for the content. An AI/LLM is not an author.

\begin{acknowledgments}
We used the FASHI DR2 and DR1 catalogues (FAST; \citealt{fashi_dr2,fashi_dr1}), the Tempel et al.\ (2017)
group catalogue, and the Douglass et al.\ (2023) SDSS void catalogue, the last two via VizieR; the analysis builds
on the \code{jansky} project. The author acknowledges the collaborative contribution of Anthropic's
Claude.
\end{acknowledgments}

\software{\code{jansky-research} \citep{janskyresearch}, Claude \citep{anthropicclaude}, astroquery, NumPy, SciPy, Matplotlib,
jansky \citep{jansky}}

\bibliography{refs}
\bibliographystyle{aasjournal}
\end{document}
